\section{Enforcement: The Two Checks}\label{sec:enforcement}

\paragraph{Check 1---tool-requirement compatibility.}
The first check enforces standard sink-side information flow control (IFC) adequacy~\cite{sabelfeld2003language}: the trust level of the trajectory must satisfy the tool's required trust floor, and the reader set must cover all intended recipients, $\readers(\fold) \supseteq \mathrm{recipients}(\mathrm{args})$, with placeholders and resolvers determining the recipient set. This mechanism adapts standard noninterference principles to tool invocation.

\paragraph{Check 2---state acquisition.}
We order $S$ by restrictiveness, where $s' \le s$ if and only if the readers of $s'$ are a subset of those in $s$ and the trust rank of $s'$ is less than or equal to that of $s$. Because every label update is a meet operation, trajectory state monotonically descends or remains fixed along both dimensions. The checker enforces the principle of least privilege at trajectory granularity~\cite{saltzer1975protection} through a prospective evaluation protocol:
\begin{enumerate}
  \item \textbf{Prospective evaluation:} The engine computes the prospective label state $s_{\mathrm{prosp}} = \fold(s, \texttt{delta(call)})$ that tool dispatch would produce.
  \item \textbf{Narrowing block:} If the call strictly restricts trajectory permissions ($s_{\mathrm{prosp}} < s$), the engine blocks execution prior to dispatch.
  \item \textbf{In-session acceptance:} Clearing the acquisition block in the active trajectory requires an explicit \texttt{Accept} step from the agent.
  \item \textbf{Atomic plan composition:} If the prospective call additionally incurs requirement-side deficits, the engine combines the \texttt{Accept} step with any necessary policy-ruling \texttt{Authorize} steps into a single atomic remedy plan.
  \item \textbf{Same-trajectory dispatch:} Upon successful execution of the remedy plan (which accepts the label narrowing and attaches any required authorizations), the tool call is dispatched within that same trajectory under the updated state $s_{\mathrm{prosp}}$.
\end{enumerate}

Trajectory forking is deliberately excluded from the mechanisms that clear an acquisition block. Because a child trajectory is seeded from the parent's current label state $s_{\mathrm{parent}}$, it inherits the same initial restrictions; consequently, spawning a child branch cannot, by itself, satisfy a state restriction or policy requirement that the parent context could not satisfy. Context branching serves as a host-level orchestration option to isolate execution under inherited labels (\S\ref{sec:branching}), whereas clearing an acquisition block within any active trajectory strictly requires explicit agent acceptance. The two verification steps remain decoupled: an external policy ruling cannot accept state narrowing on behalf of the agent, and agent acceptance cannot satisfy an unmet tool requirement.

The release frontier characterizes system-level behavior rather than individual checked objects. Because release-side requirements are monotone in $s$, state descent can only reduce the set of requests satisfiable without external rulings. Conversely, source context bounds ($\subseteq$) are antimonotone---state descent may cause previously failing bounds to be satisfied---which places them outside the scope of the release frontier theorem.

\paragraph{Label monotonicity.}
In the primary model, tool deltas only restrict access and never elevate trust or expand readership. A single tool call, such as \texttt{search\_and\_share}, may simultaneously induce a restrictive state delta and present an unmet release-side requirement. In this case, an atomic remedy plan combines an \texttt{Accept} step for state narrowing with an \texttt{Authorize} step for the requirement gap. Tools that grant access, such as \texttt{share\_doc(doc, outsider)}, are modeled as a fetch operation composed with an access-granting action, allowing each operation to be evaluated independently.

\paragraph{Evaluation timing.}
Label requirements are evaluated against the prospective state $\fold(L,\texttt{delta(call)})$ that tool dispatch would produce. This prevents race conditions where a tool call, such as \texttt{search\_and\_share} configured with \texttt{requires: audience=public} and \texttt{delta: audience=internal}, might satisfy requirements under the current state despite committing internal data upon execution. By contrast, history predicates evaluate the log in its current pre-state, ensuring that a tool call's \texttt{emits} sequence cannot satisfy its own prerequisites.

Source bounds ($\subseteq$) are prospective by default. When a data fetch restricts the label state beyond a specified bound, readers removed by the meet operation provably receive no post-fetch content. The \texttt{strict} modifier additionally enforces the requirement against the pre-state, requiring an isolated trajectory state established by an explicitly accepted prior narrowing. External communication channels that display messages to fixed readers independently of label state operations are outside the scope of this model, which governs agentic trajectories rather than arbitrary channel IFC. Memory eviction mechanisms are reserved for the permissive extension.

\begin{theorem}[Remedy completeness]\label{thm:remedy}
Relative to the registry configurations and resolver states at evaluation time, the remedy space is finite and enumerable. It comprises six distinct categories:
\begin{enumerate}
  \item \textbf{Tag-routed policy rulings} whose mandate ceilings cover the requirement gap;
  \item \textbf{Input-sanitizer substitutions} that generate valid derived arguments;
  \item \textbf{Output-sanitizer composites} in confining deployments;
  \item \textbf{Registered tools} whose \texttt{emits} sequences satisfy failed \texttt{prior(k)} predicates;
  \item \textbf{Waiver mandates and authorizing components} for demanded attention marks;
  \item \textbf{Agent acceptance} for state acquisition blocks.
\end{enumerate}
\end{theorem}
Acceptance does not require registry lookup because accepting state descent grants no security authority; consequently, acquisition blocks are non-terminal, and an empty remedy set indicates unsatisfiability only for requirement-side checks. A taint-confining branch that returns an output value requires a declared exit contract whose result label preserves the parent state; otherwise, trajectory abandonment is the sole parent-preserving outcome. For release-side failures, label monotonicity ensures completeness: unruled executions only descend state restrictiveness, and sanitizer or branch outputs merge via set intersection and minimum trust, ensuring no unenumerated operation can satisfy the gap. Thus, an empty remedy set serves as a proof of unremediability. The converse guarantee is deliberately conditional: a non-empty remedy set indicates that a valid remediation plan exists given current resolver and registry state. This property relies on enforcing that authorities with empty mandates produce an explicit initialization error (\S\ref{sec:implementation}). Proof in Appendix~\ref{app:proofs}.

\begin{proposition}[Eviction soundness, extension only]\label{prop:eviction}
Under the permissive extension, executing a restricted read following a state widening operation intersects away any reader not included in the new content's reader set. All remaining readers are entitled by construction, and evicted readers observe no post-read data.
\end{proposition}
This proposition is vacuous in the primary model, where state widening operations are absent. Proof in Appendix~\ref{app:proofs}.

\paragraph{Remedy-plan delivery.}
Remediation plans are executable objects exposed through a static tool, \texttt{execute\_remedy\_plan(plan\_id)}, initialized at run start; dynamically modifying tool definitions mid-conversation would invalidate prompt caches~\cite{kwon2023efficient}. The plan identifier, tool dispatch, and associated policy rulings are logged upon execution. Unmet attention requirements use the same blocking interface: an attention demand is treated as a failed predicate whose remediation is an atomic ruling issued by an authorizing component.

\fixme{Open question---the structure of a \texttt{prior(k)} plan: plan execution spans two calls (the event-emitting call followed by the target call); per-call verification within a single plan execution remains unresolved.}
